\section{TÁC ĐỘNG CỦA TRÍ TUỆ NHÂN TẠO (AI) ĐẾN QUẢN TRỊ VÀ QUẢN LÝ TẠI FPT}

Trong kỷ nguyên kinh tế số, Trí tuệ Nhân tạo (AI) không chỉ là một công cụ hỗ trợ mà đã trở thành trọng tâm chiến lược tái định hình tổ chức. Tại FPT, AI tạo ra tác động mạnh mẽ theo hai chiều: \textbf{vừa là mũi nhọn sản phẩm kinh doanh then chốt, vừa là công cụ nền tảng đổi mới toàn diện công tác quản trị và quản lý nội bộ}.

\subsection{AI với tư cách là sản phẩm và giải pháp kinh doanh cốt lõi}
FPT đã chủ động đón đầu làn sóng AI từ rất sớm và liên tục gặt hái những thành tựu có tính đột phá:
\begin{itemize}
    \item \textbf{Hệ sinh thái FPT.AI (khởi xướng từ năm 2017):} Cung cấp bộ giải pháp Trí tuệ Nhân tạo toàn diện phục vụ các tổ chức tài chính, ngân hàng, bảo hiểm, bán lẻ và hành chính công. Nền tảng tích hợp trợ lý ảo đàm thoại thông minh (Voicebot/Chatbot), công nghệ thị giác máy tính nhận dạng ký tự quang học (OCR) và xác thực định danh điện tử (eKYC).
    \item \textbf{Nền tảng akaBot:} Kết hợp tự động hóa quy trình bằng robot (RPA) với xử lý tài liệu thông minh (IDP) và AI đại lý (Agentic AI); triển khai thành công cho hơn \textbf{4.500 doanh nghiệp tại 21 quốc gia}, giải phóng hàng triệu giờ lao động thủ công.
    \item \textbf{Hợp tác chiến lược cùng NVIDIA xây dựng AI Factory:} Từ tháng 4/2024, FPT công bố đầu tư khoảng 200 triệu USD cùng NVIDIA xây dựng các Nhà máy Trí tuệ Nhân tạo (AI Factory) trang bị hệ thống siêu máy tính GPU Hopper tại Việt Nam và ra mắt tại Nhật Bản vào tháng 11/2024, cung cấp hạ tầng tính toán đám mây AI cho các tập đoàn toàn cầu.
    \item \textbf{Hợp đồng thương mại kỷ lục:} Tháng 9/2025, FPT chính thức ký kết hợp đồng tư vấn và triển khai chuyển đổi số ứng dụng AI có quy mô lên tới \textbf{256 triệu USD} với một tập đoàn năng lượng hàng đầu châu Á, khẳng định năng lực cạnh tranh quốc tế vượt trội.
\end{itemize}

\subsection{AI tác động đến công tác Quản trị công ty (Corporate Governance)}
Ở cấp độ quản trị chiến lược, AI đóng vai trò là ``cố vấn thông minh'' hỗ trợ HĐQT và Ban Kiểm soát:
\begin{enumerate}
    \item \textbf{Hỗ trợ ra quyết định chiến lược dựa trên dữ liệu lớn (Data-Driven Decision Making):} Các thuật toán học máy xử lý hàng triệu điểm dữ liệu thị trường, dự báo biến động kinh tế vĩ mô và phân tích tiềm năng của các đối tượng M\&A, giúp HĐQT phê duyệt các định hướng kinh doanh với độ chính xác cao.
    \item \textbf{Tự động hóa công tác giám sát và quản trị rủi ro:} Ban Kiểm soát và bộ phận Kiểm toán Nội bộ ứng dụng AI để quét tự động toàn bộ luồng giao dịch tài chính, nhanh chóng phát hiện các dấu hiệu gian lận, xung đột lợi ích hoặc vi phạm tuân thủ vượt quá hạn mức phê duyệt.
    \item \textbf{Hiện đại hóa quản trị quan hệ cổ đông (Investor Relations):} Tại Đại hội đồng Cổ đông thường niên năm 2025, hệ thống AI nhận diện khuôn mặt và tra cứu dữ liệu thời gian thực đã giúp xác thực tư cách cổ đông chỉ trong vòng \textbf{khoảng 2 giây}, nâng cao tính chuyên nghiệp, minh bạch và an toàn của đại hội.
\end{enumerate}

\subsection{AI tác động đến 4 chức năng Quản lý (Management Functions)}
Tại cấp độ quản lý điều hành, AI tác động trực tiếp làm thay đổi căn bản 4 chức năng quản lý kinh điển:

\subsubsection{1. Chức năng Hoạch định (Planning)}
AI chuyển đổi mô hình hoạch định từ ``phản ứng theo kinh nghiệm'' sang ``dự báo chủ động''. Các mô hình AI dự báo chính xác nhu cầu tuyển dụng kỹ sư công nghệ theo từng quý, dự toán ngân sách dự án và mô phỏng các kịch bản rủi ro thị trường (như biến động tỷ giá Yên Nhật hoặc Đô la Mỹ).

\subsubsection{2. Chức năng Tổ chức (Organizing)}
AI tối ưu hóa việc phân bổ nguồn lực (Resource Allocation) tại FPT Software:
\begin{itemize}
    \item Tự động hóa sàng lọc hàng trăm nghìn hồ sơ ứng viên (CV Screening), khớp nối năng lực của kỹ sư với yêu cầu công nghệ của từng dự án phần mềm cụ thể.
    \item Dự báo tỷ lệ nghỉ việc (attrition rate) của nhân viên dựa trên các chỉ số gắn kết để nhà quản lý chủ động can thiệp giữ chân nhân tài.
\end{itemize}

\subsubsection{3. Chức năng Lãnh đạo (Leading)}
Hệ thống phân tích thông minh cung cấp cho các nhà quản lý bức tranh toàn cảnh về tâm lý và hiệu suất làm việc của nhóm. Các chỉ số tương tác nội bộ, phản hồi ẩn danh được AI phân tích giúp lãnh đạo cải tiến phương thức truyền cảm hứng, xây dựng môi trường làm việc hạnh phúc.

\subsubsection{4. Chức năng Kiểm soát (Controlling)}
FPT triển khai các bảng điều khiển thời gian thực (Real-time Executive Dashboards) tích hợp AI, cho phép CEO và các giám đốc khối giám sát tiến độ dự án, năng suất lao động và sai sót kỹ thuật theo từng giây mà không cần chờ báo cáo định kỳ bằng văn bản giấy tờ.

\subsubsection{Đột phá trong hiệu suất sản xuất phần mềm}
Tại FPT Software, việc tích hợp các trợ lý lập trình AI (AI Coding Assistants như GitHub Copilot hoặc giải pháp Made-by-FPT) đã giúp tự động hóa việc sinh mã nguồn, phát hiện lỗ hổng bảo mật và tự động tạo ca kiểm thử (Unit Test), giúp \textbf{tăng năng suất lao động của đội ngũ kỹ sư từ 20\% đến 30\%}.

\subsection{Thách thức và rủi ro quản trị khi ứng dụng AI}
Bên cạnh những lợi ích to lớn, việc ứng dụng AI đặt ra những thách thức sống còn đối với ban lãnh đạo FPT:
\begin{itemize}
    \item \textbf{Rủi ro an ninh mạng và bảo mật dữ liệu khách hàng:} Việc các kỹ sư sử dụng mô hình AI tạo sinh (GenAI) tiềm ẩn nguy cơ rò rỉ mã nguồn độc quyền của khách hàng toàn cầu, vi phạm nghiêm trọng các cam kết bảo mật (NDA).
    \item \textbf{Vấn đề đạo đức và tính giải trình của AI (AI Ethics \& Explainability):} Các thuật toán ra quyết định tự động có thể mang thiên kiến tiềm ẩn trong đánh giá năng lực nhân viên hoặc lựa chọn nhà cung cấp, đòi hỏi khung pháp lý nội bộ minh bạch.
    \item \textbf{Áp lực đào tạo lại kỹ năng (Reskilling):} Đòi hỏi FPT phải đầu tư ngân sách lớn để đào tạo phổ cập kiến thức AI cho đội ngũ gần 90.000 nhân sự trong toàn hệ sinh thái, tránh nguy cơ tụt hậu năng suất.
\end{itemize}
